\documentclass[a1paper,portrait]{baposter}



\usepackage{wrapfig}
\usepackage{lmodern}

\usepackage[utf8]{inputenc} %unicode support
\usepackage[T1]{fontenc}
\usepackage{float}


\selectcolormodel{cmyk}

\usepackage[skip=2pt,font=footnotesize]{caption}

\usepackage{enumitem}
\usepackage{natbib}
\usetikzlibrary{backgrounds}

\usepackage{amsmath}
\usepackage{caption}
\usepackage{subcaption}
\usepackage{url}
\usepackage{wrapfig}


\setlist{nolistsep}

\graphicspath{{figures/}} % Directory in which figures are stored

\usepackage{lipsum}

\newcommand{\compresslist}{%
\setlength{\itemsep}{0pt}%
\setlength{\parskip}{1pt}%
\setlength{\parsep}{0pt}%
}

\newenvironment{boenumerate}
  {\begin{enumerate}\renewcommand\labelenumi{\textbf\theenumi.}}
  {\end{enumerate}}


\usepackage{multirow}
\usepackage{algorithm}
\usepackage[noend]{algpseudocode}
\usepackage{algorithmicx}
%\usepackage{algorithmic}

\begin{document}
\graphicspath{ {images/} }


%%===========%%
%% Uni Wien colours %%
%%===========%%

\definecolor{uniwienblack}{RGB}{0,0,0}
\definecolor{uniwienwhite}{RGB}{255,255,255}
\definecolor{uniwienblue}{RGB}{0,99,166}
\definecolor{uniwienred}{RGB}{167,28,73}
\definecolor{uniwienorange}{RGB}{221,72,20}
\definecolor{uniwienyellow}{RGB}{246,168,0}
\definecolor{uniwiengreen}{RGB}{148,193,84}
\definecolor{uniwiencyan}{RGB}{17,137,122}

\background{ % Set the background to an image (background.pdf)
	\begin{tikzpicture}[remember picture,overlay]
	\draw (current page.south west)+(-.65em,-.5em) node[anchor=south west]
	\end{tikzpicture}
}


\begin{poster}
{
grid=False,
headerborder=open, % Adds a border around the header of content boxes
colspacing=1em, % Column spacing
bgColorOne=white, % Background color for the gradient on the left side of the poster
bgColorTwo=white, % Background color for the gradient on the right side of the poster
borderColor=uniwienblue, % Border color
headerColorOne=uniwienblue, % Background color for the header in the content boxes (left side)
headerColorTwo=uniwienblue, % Background color for the header in the content boxes (right side)
headerFontColor=white, % Text color for the header text in the content boxes
boxColorOne=white, % Background color of the content boxes
textborder=rounded, %rectangle, % Format of the border around content boxes, can be: none, bars, coils, triangles, rectangle, rounded, roundedsmall, roundedright or faded
eyecatcher=false, % Set to false for ignoring the left logo in the title and move the title left
headerheight=0.11\textheight, % Height of the header
headershape=rounded, % Specify the rounded corner in the content box headers, can be: rectangle, small-rounded, roundedright, roundedleft or rounded
headershade=plain,
headerfont=\Large\textsf, % Large, bold and sans serif font in the headers of content boxes
%textfont={\setlength{\parindent}{1.5em}}, % Uncomment for paragraph indentation
linewidth=2pt, % Width of the border lines around content boxes
%background=user
columns=4
}
{}
%
%----------------------------------------------------------------------------------------
%	TITLE AND AUTHOR NAME
%----------------------------------------------------------------------------------------
%
{
\textsf %Sans Serif
{Put your amazing poster title here}
}
{
    \sf\vspace{+0.65em}\\
    First Author Name, Second Author Name\\
    \small{
        Department of Business Decisions and Analytics, University of Vienna, 1090 Vienna, Austria \\
        {first.last@univie.ac.at}, first.last@univie.ac.at
        (This research is funded through XX project No. YYY) % remove it not applicable
    } \\
    \vspace{-1.5em}
}
% University/lab logo
{
    \includegraphics[width=0.25\textwidth]{images/logo_uni_wien.jpg}
} % ignore the error here, we need the brackets for correct formatting


%%% --- Top part --- %%%

% abstract
\headerbox{Abstract}{name=abstract,column=0,row=0, span=4}{
\lipsum[1]
}

%%% --- START middle part --- %%%

%% left column

% skill vrp
\headerbox{Section 1}{name=sec1, column=0, below=abstract, span=2}{

\lipsum[2-6]

}

%% right column

% deep reinforcement learning
\headerbox{Section 2}{name=sec2, column=2, span=2, below=abstract}{

\lipsum[7-8]

}

% training
\headerbox{Section 3}{name=sec3, column=2, span=2, below=sec2}{

\lipsum[9-10]

}

%%% --- END middle part --- %%%

%%% --- Bottom part --- %%%

% results
\headerbox{Section 4/Results}{name=model, span=3, column=0, below=sec1, above=bottom}{

\lipsum[11-12]

}

% references
\headerbox{References}{name=references,column=3,span=1,below=sec1, above=bottom}{ % for "below" reference whichever box reaches the furthest down.
\setlength{\bibsep}{2pt plus 0.3ex}
%\fontsize{5pt}{5pt}
\tiny % Reduce the font size in this block
%\scriptsize
\renewcommand{\section}[2]{\vskip 0.05em} % Get rid of the default "References" section title
\nocite{*} % Insert publications even if they are not cited in the poster


\bibliographystyle{unsrt}
\bibliography{bib.bib} % Use sample.bib as the bibliography file
%\vspace{-1mm}
}
\end{poster}

\end{document}
